\section{Several lost divisor directions}\label{sec:multiple}

We now retain the geometric hypotheses of Section~\ref{sec:introduction}
but replace \eqref{eq:betti-hypotheses} by
$b_2(Y)=b$ and $b_2(X_t)=s$, with $m=b-s>0$.
The local metric, scalar and capacity arguments above do not depend on
these dimensions. The additional issue is global: the trial forms must
be projected off the entire $s$-dimensional harmonic space, and a lower
bound is needed on every subspace of the small spectral cluster.

\subsection{The full harmonic space and its quotient norm}
Put $V=H^2(Y;\RR)$ and let $r:V\to\RR^k$ be the exceptional-period
map. The Thom and Mayer--Vietoris arguments of
Proposition~\ref{prop:restriction-spectrum} give
\[
 K:=\ker r=\im\bigl(H^2(X_t;\RR)\longrightarrow V\bigr),
 \qquad \dim K=s.
\]
In particular $H\in K$. Write
\begin{equation}\label{eq:full-intersection-forms}
 B(D,E)=\int_YDEH,\qquad a(D)=\int_YDH^2,\qquad v=\int_YH^3.
\end{equation}
Simple connectivity transfers to the smoothing by van Kampen, as before.
Both smooth spaces have $H^{2,0}=0$, so integral second-cohomology
classes are represented by line bundles.

Apply the curvature construction and scalar correction in the proof of
Proposition~\ref{prop:mass} to an integral basis of $V$, and extend
linearly. For $D\in V$ this gives a closed and coclosed real $(1,1)$-form
$\eta_D$ with class $D$, constant contraction $3a(D)/v$, and link
coefficients $r_i(D)$. The scalar inverse is fixed and linear; its
solvability condition uses $a(D)$. The local Picard calculation,
transgression estimates and logarithmic domain cutoffs do not use the
dimension of $V$. After taking the minimum of finitely many positive
rates, the resulting estimates hold uniformly on bounded subsets of $V$.
Polarizing \eqref{eq:mass-identity} gives
\begin{equation}\label{eq:full-harmonic-gram}
 G(D,E):=\langle\eta_D,\eta_E\rangle
 =-B(D,E)+\frac{3a(D)a(E)}{2v}.
\end{equation}
Indeed the primitive parts have pairing
$-B(D,E)+a(D)a(E)/v$, and $\|\omega\|^2=v/2$.
The independent classes of the $\eta_D$ make $G$ positive definite.

We will need uniqueness for limits on the actual central metric, as
well as on the exact-cone modifications.

\begin{lemma}\label{lem:actual-harmonic-uniqueness}
Let a smooth metric on $M_0$ be uniformly equivalent near each end to
the six-dimensional cone metric. Every smooth globally exact $L^2$
two-form $u$ has an $L^2$ primitive in $\operatorname{Dom}d_{\min}$.
Consequently a smooth $L^2$ closed and formally coclosed two-form is
determined by its ordinary de Rham class.
\end{lemma}
\begin{proof}
On one end write $u=a(r)+dr\wedge b(r)$ in fixed link coordinates.
Choose $v_0$ with $d_Lv_0=a(r_0)$ and set
\[
 \beta(r)=v_0+\int_{r_0}^r b(x)\,dx.
\]
Closedness gives $d\beta=u$. In the fixed link norm put
$B_0(x)=\|b(x)\|_{L^2(L,h)}$. Cone equivalence then implies
$\int_0^{r_0}x^3B_0(x)^2dx<\infty$.
For $0<r<\delta<r_0$, Cauchy--Schwarz gives
\[
 r\|\beta(r)\|_{L^2(L)}
 \le rC_\delta+2^{-1/2}
       \left(\int_0^\delta x^3B_0(x)^2dx\right)^{1/2}.
\]
First let $r\to0$, then $\delta\to0$. Thus
$r\|\beta(r)\|_{L^2(L)}\to0$. The one-form $\beta$ is $L^2$,
since its squared cone norm has radial weight $r^3$. If $\chi_\rho$
vanishes below $\rho$ and equals one above $2\rho$, with derivative
bounded by $C/\rho$, then
\[
 \|d\chi_\rho\wedge\beta\|^2
 \le C\rho^{-2}\int_\rho^{2\rho}
        r^3\|\beta(r)\|_{L^2(L)}^2dr=o(1).
\]
The $L^2$ tails of $\beta$ and $d\beta=u$ also vanish.

Let $w_0$ be a global smooth primitive. On the end $w_0-\beta$ is
closed and hence exact, since $H^1(L)=0$. Subtracting the differential
of a cutoff potential patches $w_0$ to $\beta$ near each tip without
changing its differential. The resulting global $w$ is $L^2$, and the
preceding cutoff estimate proves $w\in\operatorname{Dom}d_{\min}$.
If $u$ is formally coclosed, testing on the compactly supported core
and passing to its graph closure gives
$u\in\operatorname{Dom}d_{\min}^*$ and $d_{\min}^*u=0$. Hence
$\|u\|^2=\langle d_{\min}w,u\rangle=0$.
Apply this to the difference of two forms to obtain uniqueness.
\end{proof}

For the actual central metric, this lemma and the $b$ constructed forms
identify the smooth $L^2$ harmonic two-form space with $V$.
For $g_0^R$, the proof of Lemma~\ref{lem:central-kernels} applies to
a basis of $b$ classes: compactly supported dual four-forms give $b$
independent reduced classes, while its local expansion gives injection
into ordinary cohomology. Its degree-one kernel is still zero.
Denote the resulting central harmonic Gram form by $G_R$.
Minimization in the same closed reduced complex and norm equivalence
give $G_R\to G$, exactly as in \eqref{eq:mass-limit}.

Give $\mathcal Q=V/K$ the quotient norm $M$ induced by $G$.
The $G$-orthogonal lift $D^\perp$ of $[D]$ is primitive, because
$G(D,H)=a(D)/2$. Thus
\begin{equation}\label{eq:quotient-mass-matrix}
 M([D],[E])=-B(D^\perp,E^\perp),\qquad
 D^\perp\perp_B K.
\end{equation}
The restriction of $B$ to $K$ is nondegenerate: its decomposition
$K=\RR H\oplus(K\cap\ker a)$ has one positive direction and a
negative-definite primitive part. In a basis of $K$ followed by
quotient lifts, the mass matrix is therefore
\begin{equation}\label{eq:mass-schur}
 \mathsf M=-B_{DD}+B_{DK}B_{KK}^{-1}B_{KD}.
\end{equation}
The period form
$E([D],[E])=r(D)\cdot r(E)$ is positive definite on $\mathcal Q$.
For $g_0^R$ let $M_R$ be the quotient norm induced by $G_R$.
Then $M_R\to M$, and the generalized eigenvalues
$\mu_j(R)$ of $(E,M_R)$ converge to those of $(E,M)$.
Subtracting the polarization component alone would not give this
quotient mass when $s>1$.

\subsection{The surviving harmonic directions}
For $k\in K$, the form $\eta_k$ has zero harmonic link coefficient.
It admits on each end a primitive $\beta_k$ of cone size
$O(r^{\nu-1})$, with differentiated bounds. Indeed, write
$\eta_k=a(r)+dr\wedge b(r)$. Its tangential coefficients in fixed
link coordinates are $O(r^\nu)$, and
$\partial_ra=d_Lb$. Thus $\beta_k(r)=\int_0^r b(x)dx$ is a
primitive with the claimed estimate. The bounded K\"ahler component
satisfies the same bound after decreasing $\nu$.

Choose $\chi_\rho$ equal one near the tip and zero for $r\ge2\rho$.
Replacing $\eta_k$ on each end by
\begin{equation}\label{eq:closed-compact-approximation}
 \gamma_{k,\rho}=\eta_k-d(\chi_\rho\beta_k)
\end{equation}
gives a compactly supported closed form representing $k$, converging
in $L^2$ to $\eta_k$. The removed tail and cutoff correction have
squared norm $O(\rho^{2\nu+2})$.
For fixed $g_0^R$ its exact-cone expansion gives the same construction.

Let $h_t(k)$ be the actual harmonic representative of the corresponding
class on $X_t$. Transport \eqref{eq:closed-compact-approximation} and
extend it by zero. It remains closed in that class, so minimization gives
\[
 \limsup_{t\to0}\|h_t(k)\|^2
 \le\|\gamma_{k,\rho}\|_{g_0}^2.
\]
Letting $\rho\to0$ bounds this by $G(k,k)$. Local elliptic estimates
give a subsequential smooth formally harmonic limit $h_0$ on $M_0$,
which is $L^2$ by Fatou. Pairing with compactly supported closed dual
four-forms identifies its class as $k$. Lemma~\ref{lem:actual-harmonic-uniqueness}
applied to $h_0-\eta_k$ therefore gives $h_0=\eta_k$.
Lower semicontinuity proves the reverse norm inequality.
Consequently the norms converge, no mass remains in the shrinking ends,
and the convergence holds for the full family. Polarization over a
basis proves convergence of the whole harmonic Gram matrix on $K$.
The same proof works for each fixed exact-cone model with $G_R$.

\subsection{Upper bounds on all small subspaces}
For $D\in K^{\perp_G}$ apply the logarithmic trial construction of
Section~\ref{sec:capacity} linearly to $\eta_D$. Write $T_\epsilon D$
for the transported form, $L_\epsilon=\log(1/\epsilon)$ and $c=2\pi/3$.
The estimates there, polarized on this fixed finite-dimensional space,
give
\begin{align}
 q_\epsilon(T_\epsilon D)
  &=\frac{cE(D,D)}{L_\epsilon}
          +O(L_\epsilon^{-2})\|D\|^2,\label{eq:matrix-trial-energy}\\
 \langle T_\epsilon D,T_\epsilon E\rangle
  &=M(D,E)+O(L_\epsilon^{-1})\|D\|\|E\|.\label{eq:matrix-trial-mass}
\end{align}
The finite-dimensional norm on the right is arbitrary and fixed.
The same mass estimates give strong exterior convergence to $\eta_D$
and uniformly vanishing $L^2$ tails on bounded sets of $D$.

Project orthogonally onto actual real $(1,1)$-forms. The K\"ahler
Hodge Laplacian preserves bidegree, so this projection does not increase
$q_\epsilon$. It changes the $L^2$ norm by $o(1)$: on fixed exterior
compacts the complex structures converge and $\eta_D$ is of type
$(1,1)$; the tail estimate controls the remaining region. Next subtract
the full harmonic projection. Its norm is $o(1)$ by the preceding
subsection and $G(D,K)=0$, and its removal leaves the energy unchanged.
The resulting real $(1,1)$-forms $v_\epsilon(D)$ are orthogonal to
the harmonic kernel, with mass Gram form tending to $M$.

Let $\mu_1\le\cdots\le\mu_m$ be the generalized eigenvalues of
$(E,M)$. Apply min--max to the image under $v_\epsilon$ of the first
$j$ generalized eigendirections. Its dimension is $j$ for small
$\epsilon$, and \eqref{eq:matrix-trial-energy} gives
\begin{equation}\label{eq:matrix-upper}
 \limsup_{\epsilon\to0}L_\epsilon\lambda_{2,j}^+(\epsilon)
 \le c\mu_j,\qquad 1\le j\le m.
\end{equation}
The same bound holds within real $(1,1)$-forms. Only an $o(1)$ estimate
was needed for the full kernel projection; this argument does not claim
the stronger upper remainder of the rank-one theorem.

\subsection{Counting and the ordered lower bounds}
Fix the auxiliary radii $R,A$ in the exact-cone model. The cap-kernel
vanishing and physical-degree argument in
Proposition~\ref{prop:model-count} remain unchanged. The norm
decomposition preserves every finite orthonormal collection of
zero-limit one- or two-form eigenvectors as a central collection,
without cap or critical mass. The central dimensions are now zero
and $b$. Conversely graph-norm approximation of a basis of all $b$
central harmonic two-forms gives $b$ zero-limit eigenvalues by
min--max. Exactly $s$ of these are the topological zeros. Hence there
are precisely $m$ positive small eigenvalues, and the complementary
degree-one and positive degree-two gaps hold. Their orthonormal central
limits belong to $K^{\perp_{G_R}}$, because the model's full harmonic
kernel converges to $K$ by the preceding argument.

Write $\bar\lambda_j(\epsilon;R,A)$ for these ordered positive
model eigenvalues. Choose a sequence on which
$L_\epsilon\bar\lambda_j$ has finite lower limit $\ell$; an
infinite lower limit requires no proof. The first $j$ orthonormal
eigenforms have a subsequence with orthonormal central limits
$v_1,\ldots,v_j\in K^{\perp_{G_R}}$. Every fixed linear combination
$u_\epsilon(z)$ satisfies
\[
 q(u_\epsilon(z))\le\bar\lambda_j|z|^2,
 \qquad v(z)=\sum_{a=1}^jz_av_a.
\]
At the fixed sections $r_*=R/2$, its harmonic boundary coefficients
converge to $r_i(v(z))$. Apply Lemma~\ref{lem:capacity} to this
combination on every disjoint cap region and sum. It follows that
\[
 cE(v(z),v(z))\le\ell |z|^2
                  =\ell M_R(v(z),v(z)).
\]
One subsequence suffices for the finite basis, hence for every fixed
combination; equivalently first use a countable dense set of $z$ and
then continuity. The largest generalized Rayleigh quotient on this
$j$-dimensional span is at least $\mu_j(R)$. Therefore
\begin{equation}\label{eq:matrix-model-lower}
 \liminf_{\epsilon\to0}L_\epsilon\bar\lambda_j(\epsilon;R,A)
 \ge c\mu_j(R).
\end{equation}

The derivative-free comparison transfers this inequality to the actual
metrics. Taking the family limit first, then $A\to\infty$ and
$R\to0$, removes its multiplicative error and sends $\mu_j(R)$ to
$\mu_j$. Together with \eqref{eq:matrix-upper} and
$L_\epsilon=|\log t|/3$, this proves the law in
Theorem~\ref{thm:matrix}. One fixed admissible pair $R,A$ transfers
the complementary positive gaps.

Finally, \eqref{eq:matrix-upper} supplies $m$ small eigenvalues in
real $(1,1)$-forms, and the total count is exactly $m$. Thus this type
exhausts the small cluster. The degree-one gap excludes an exact small
two-form. Contraction by $\Lambda_\omega$ commutes with the K\"ahler
Laplacian; a nonzero contraction would give a positive scalar eigenfunction
whose differential contradicts that same gap. Hence every small
eigenform is primitive and coexact. A simple generalized eigenvalue
$\mu_j$ is isolated from its neighbours, so the leading law and total
count imply eventual simplicity. This completes the proof.
